\documentclass{article}

% Recommended, but optional, packages for figures and better typesetting:
\usepackage{microtype}
\usepackage{graphicx}
\usepackage{subcaption}
\usepackage{booktabs} % for professional tables
\usepackage{hyperref}

% Attempt to make hyperref and algorithmic work together better:
\newcommand{\theHalgorithm}{\arabic{algorithm}}

\usepackage{icml2026}

\usepackage{amsmath}
\usepackage{amssymb}
\usepackage{mathtools}
\usepackage{amsthm}

% if you use cleveref..
\usepackage[capitalize,noabbrev]{cleveref}

%%%%%%%%%%%%%%%%%%%%%%%%%%%%%%%%
% THEOREMS
%%%%%%%%%%%%%%%%%%%%%%%%%%%%%%%%
\theoremstyle{plain}
\newtheorem{theorem}{Theorem}[section]
\newtheorem{proposition}[theorem]{Proposition}
\newtheorem{lemma}[theorem]{Lemma}
\newtheorem{corollary}[theorem]{Corollary}
\theoremstyle{definition}
\newtheorem{definition}[theorem]{Definition}
\newtheorem{assumption}[theorem]{Assumption}
\theoremstyle{remark}
\newtheorem{remark}[theorem]{Remark}

\icmltitlerunning{Quantization-Constrained Sliced Wasserstein}

\begin{document}

\twocolumn[
\icmltitle{Quantization-Constrained Sliced Wasserstein:\\
           Multi-Dimensional Optimal Transport for Robust Model Merging}

\begin{icmlauthorlist}
\icmlauthor{The Theorist}{inst}
\end{icmlauthorlist}

\icmlaffiliation{inst}{Department of Mathematical Foundations in AI, Switzerland}
\icmlcorrespondingauthor{The Theorist}{theorist@iaml.ch}

\icmlkeywords{Model Merging, Sliced Wasserstein, Optimal Transport, Post-Training Quantization}

\vskip 0.3in
]

\printAffiliationsAndNotice{}

\begin{abstract}
Post-hoc multi-task model merging has emerged as a promising paradigm to consolidate specialized expert neural networks into a single cohesive model. However, direct linear averaging of weights triggers catastrophic representation collapse due to high-dimensional update orthogonality, which is exacerbated under low-bit post-training quantization (PTQ). While recent 1D Wasserstein calibration (QCOT) prevents collapse by parameter-space alignment under infinity-norm constraints, it independently aligns 1D coordinate marginals, ignoring rich multi-dimensional joint feature distributions of weight vectors (neurons) and introducing geometric distortions on complex datasets. In this paper, we propose Quantization-Constrained Sliced Wasserstein (QCSW) calibration, a mathematically rigorous multi-dimensional optimal transport framework for robust model merging. QCSW projects multi-dimensional neuron weight vectors onto random 1D lines, performs optimal quantization-constrained transport in the sliced spaces, and reconstructs the calibrated high-dimensional parameters via closed-form pseudo-inverse. We prove that QCSW converges to a multi-dimensional Wasserstein projection under update constraints. Empirically, QCSW achieves outstanding multitask accuracy on ResNet-18 across MNIST, Fashion-MNIST, and CIFAR-10 tasks under severe PTQ and environmental noise, outperforming standard WCPR, QCOT, and task arithmetic baselines.
\end{abstract}

\section{Introduction}
\label{sec:intro}

The pre-train and fine-tune paradigm has enabled the construction of specialized task-specific expert neural networks from a single, shared progenitor backbone~\cite{devlin2018bert, he2016deep, vaswani2017attention}. However, deploying and serving multiple independent backbones for diverse downstream tasks incurs prohibitive memory, storage, and scheduling overheads, especially on edge hardware. To bypass this server-serving bottleneck, multi-task model merging has emerged as a practical, training-free approach to consolidate independent expert networks into a single, cohesive multi-task model without parameter training or access to the massive original datasets~\cite{wortsman2022model, ilharco2022editing}.

Despite its high cost-efficiency, simple uncalibrated parameter combination, such as Weight Averaging (WA) or Task Arithmetic (TA)~\cite{ilharco2022editing}, triggers severe performance degradation in deep layers~\cite{jordan2023repair}. This pathology, known as ``representation collapse'' or ``variance decay,'' occurs because fine-tuned task updates (task vectors) are approximately orthogonal in high-dimensional parameter space. Summing or averaging these updates reduces the Frobenius norm of deep layers by a factor of $\sqrt{K}$ (where $K$ is the number of experts), which compounds exponentially across layers and results in highly distorted hidden activations.

To heal representation collapse, data-free and offline parameter-space calibration methods have been proposed. Notable among these is Wasserstein-Calibrated Parameter Resonance (WCPR)~\cite{leclaire2024wcpr}, which uses 1D Optimal Transport (OT) to non-parametrically align the empirical merged weight distributions channel-by-channel with the exact Wasserstein-2 barycenter of the experts. While WCPR perfectly restores higher-order statistical moments, it is highly vulnerable to Post-Training Quantization (PTQ) and environmental noise due to unconstrained dynamic range inflation. To resolve this, Quantization-Constrained Optimal Transport (QCOT) formulates parameter calibration as a constrained optimal transport problem, proving that the unique optimal monotonic mapping corresponds to a clipped 1D Wasserstein barycenter projection.

However, both WCPR and QCOT are strictly limited to 1D optimal transport. They treat individual parameter coordinates within each channel as independent, identically distributed (i.e.d.) scalar samples, ignoring the rich multi-dimensional joint feature structures of weight vectors (such as neuron rows in a linear layer or filters in a convolutional layer). Projecting and aligning these 1D marginals separately introduces weight-space distortions that degrade features on complex datasets (such as CIFAR-10) under low-bit post-training quantization (PTQ). 

In this work, we present **Quantization-Constrained Sliced Wasserstein (QCSW)** calibration, the first mathematically rigorous multi-dimensional optimal transport theory for robust, zero-shot model merging. We treat each neuron's weights as multi-dimensional vectors and formulate parameter alignment as minimizing the Sliced Wasserstein distance (SWD) between the merged parameter distribution and the expert barycenter, subject to strict infinity-norm quantization constraints. Our contributions are:

1. **Theoretical Formulation:** We present the first multi-dimensional, quantization-constrained optimal transport formulation for model merging using Sliced Wasserstein Geometry.
2. **Closed-Form Solver:** We prove that the infinite-dimensional optimization problem along each slice has a unique, exact closed-form solution. We reconstruct the high-dimensional calibrated parameters via a closed-form pseudo-inverse of the projection matrix.
3. **Mathematical Guarantees:** We provide formal theorems proving that QCSW calibration is unbiased (preserves the first moment) and strictly bounds the variance distortion under quantization constraints.
4. **Empirical Success:** Empirically, QCSW achieves outstanding multitask accuracy on ResNet-18 across MNIST, Fashion-MNIST, and CIFAR-10 tasks under severe PTQ and environmental noise, outperforming standard WCPR, QCOT, and task arithmetic baselines.

\section{Related Work}
\label{sec:related}

Our work intersects three active domains of machine learning theory: model merging, post-training quantization, and optimal transport.

\subsection{Model Merging and Weights Fusion}
Consolidating decentralized expert networks has emerged as a crucial area of research. Weight Averaging (WA) and Model Soups~\cite{wortsman2022model} average fine-tuned models sharing a backbone. Task Arithmetic (TA)~\cite{ilharco2022editing} views fine-tuning as adding steering vectors. To address permutation symmetries, Git Re-Basin~\cite{ainsworth2022git} and related permutation-invariant connection studies~\cite{frankle2020linear, entezari2021role} align neurons via Hungarian matching. Fisher Weighted Averaging~\cite{matena2022merging} merges parameters using diagonal Fisher information matrices. ZipIt!~\cite{stoica2023zipit} merges features of models trained on disjoint tasks, while TIES-Merging~\cite{yadav2023resolving} prunes and resolves interference. DARE~\cite{jin2023dare} employs random pruning to sparsify expert updates, and REPAIR~\cite{jordan2023repair} resolves activation scale collapse by scaling weight magnitudes. Other approaches include language model super-mergers~\cite{yu2023language}, elastic merging~\cite{daheim2023elastic}, and multi-task fusion~\cite{choshen2022fusing, ortiz2023task, donyehiya2022cold, gupta2020model, yang2020fusing}. However, these algorithms do not explicitly protect multi-task networks from representation collapse under post-training quantization constraints.

\subsection{Post-Training Quantization (PTQ)}
PTQ compresses model parameters without training data~\cite{gholami2021survey}. Advanced methods like AdaRound~\cite{nagel2020up} optimize weight rounding via local objectives. GPTQ~\cite{frantar2022gptq} and QLoRA~\cite{dettmers2023qlora} quantize weight matrices to 4-bit precision via second-order Taylor approximations. SmoothQuant~\cite{xiao2023smoothquant} and AWQ~\cite{lin2023awq} mitigate activation outliers by scaling input channels. Data-free weight equalization and bias correction~\cite{nagel2019data} reduce quantization error. Classic works explore low-bit bounds and scaling factors~\cite{banner2019post, zhao2019improving, hubara2016quantized, jacob2018quantization, krishnamoorthi2018quantizing, esser2019learned, choi2018pact, bhalgat2020lsq}. While these techniques optimize single-model compression, they often fail under multi-task model merging, where weight dynamics are disrupted.

\subsection{Optimal Transport and Sliced Wasserstein Distance}
Optimal Transport (OT) defines rigorous metrics between probability measures~\cite{kantorovich1942translocation, monge1781memoire, wasserstein1969markov, brenier1991polar, peyre2019computational, villani2009optimal}. OT methods are widely used in machine learning, such as in Wasserstein GANs~\cite{arjovsky2017wasserstein}, supervised learning~\cite{frogner2015learning}, Sinkhorn distances~\cite{cuturi2013sinkhorn}, and general domain adaptation~\cite{kolouri2017optimal, genevay2018learning}. Sliced Wasserstein Distance (SWD) solves high-dimensional OT by projecting distributions onto 1D lines where closed-form solutions exist~\cite{rabin2011wasserstein, bonneel2015sliced}. SWD has been applied in generative modeling~\cite{deshpande2018generative}, autoencoders~\cite{meng2019sliced}, and multinomial learning~\cite{kolouri2019sliced, nadjahi2019asymptotic, nadjahi2020approximate}. Hierarchical and kernel variations further expand its applicability~\cite{nguyen2021hierarchical, carriere2020sliced, bonet2021sliced}. In parameter fusion, OT has aligned neuron weights across models~\cite{singh2020model}, and WCPR~\cite{leclaire2024wcpr} aligned 1D marginals with the Wasserstein barycenter. Our work is the first to integrate Sliced Wasserstein geometry with quantization-constrained optimal transport for multi-dimensional model merging.

\section{Theoretical Framework}
\label{sec:theory}

\subsection{Preliminaries and Representation Collapse}

Let $W^{(l)}_{\text{init}} \in \mathbb{R}^{D_{\text{out}} \times D_{\text{in}}}$ be the weight matrix of a shared pre-trained progenitor model at layer $l$. Let $W^{(l)}_k = W^{(l)}_{\text{init}} + T^{(l)}_k$ be the parameters of expert $k \in \{1, \dots, K\}$, where $T^{(l)}_k$ is the task vector update. Under Weight Averaging (WA), the merged parameters are defined as:
\begin{equation}
W^{(l)}_{\text{merged}} = W^{(l)}_{\text{init}} + T^{(l)}_{\text{merged}}, \quad T^{(l)}_{\text{merged}} = \frac{1}{K} \sum_{k=1}^K T^{(l)}_k
\end{equation}

Since expert updates are highly orthogonal in deep layers ($\langle T^{(l)}_i, T^{(l)}_j \rangle_F \approx 0$ for $i \neq j$), the Frobenius norm of the merged update decays by a factor of $\sqrt{K}$:
\begin{equation}
\|T^{(l)}_{\text{merged}}\|_F \approx \frac{\|T^{(l)}\|_F}{\sqrt{K}}
\end{equation}
This norm collapse compounds exponentially across layers, causing hidden activation scales to shrink exponentially, triggering representation collapse~\cite{jordan2023repair}.

\subsection{1D Wasserstein Calibration (WCPR \& QCOT)}

To restore the collapsed statistics of the merged updates without using real calibration data, WCPR models each task update channel-by-channel as an empirical probability distribution. In 1D, the unique optimal transport map $T : \mathbb{R} \to \mathbb{R}$ pushing the empirical distribution of $T^{(l)}_{\text{merged}}[c]$ ($\mu$) to the exact Wasserstein-2 barycenter of the experts ($\nu$) has a closed-form solution based on sorting. For empirical arrays of equal size $N$, let $x_{(1)} \le \dots \le x_{(N)}$ be the sorted elements of the merged task vector channel $T^{(l)}_{\text{merged}}[c]$, and let $s_{k, (1)} \le \dots \le s_{k, (N)}$ be the sorted task vector updates of expert $k$. WCPR maps the merged parameters to the arithmetic average of the sorted experts:
\begin{equation}
T^*(x_{(i)}) = y_{(i)} := \frac{1}{K} \sum_{k=1}^K s_{k, (i)}
\end{equation}

While WCPR perfectly restores all moments, it is unconstrained and vulnerable to outlier elements $y_{(1)}$ and $y_{(N)}$, which inflate the dynamic range. QCOT resolves this by formulating a constrained optimal transport problem:
\begin{equation}
\min_{T \in \mathcal{M}} \mathcal{W}_2^2(T_\# P_X, P_Y) \quad \text{s.t.} \quad \|T(X)\|_\infty \le C
\end{equation}
where $C > 0$ is a robust, quantization-stable clipping threshold. The unique optimal monotonic mapping has the closed-form expression:
\begin{equation}
T^*(x_{(i)}) = \text{clip}(y_{(i)}, -C, C)
\end{equation}

\section{Proposed Method: Quantization-Constrained Sliced Wasserstein (QCSW)}
\label{sec:method}

\subsection{Limitation of 1D Coordinate Alignment}
The core limitation of WCPR and QCOT is that they treat each coordinate within a channel independently, ignoring the joint distribution of the weight vectors. In a neural network layer, each of the $D_{\text{out}}$ rows in the weight matrix represents a neuron's multi-dimensional weight vector in $\mathbb{R}^{D_{\text{in}}}$. Treating these as independent 1D coordinates disrupts the rich multi-dimensional geometric feature correlation structures, causing severe feature distortions on complex tasks (e.g., CIFAR-10) under quantization.

\subsection{The QCSW Formulation}
To resolve this, we propose aligning the joint multi-dimensional neuron distributions across experts using **Sliced Wasserstein Distance (SWD)**. The Sliced Wasserstein distance computes the distance between multi-dimensional distributions by projecting them onto 1D lines, solving the closed-form 1D OT problem, and averaging over all projection directions.

We formulate the Sliced Wasserstein calibration as finding a multi-dimensional calibrated update matrix $\tilde{T} \in \mathbb{R}^{D_{\text{out}} \times D_{\text{in}}}$ that minimizes the expected Wasserstein-2 distance of the projections onto the unit sphere $\mathbb{S}^{D_{\text{in}}-1}$, subject to an update infinity-norm constraint:
\begin{equation}
\min_{\tilde{T}} \mathbb{E}_{\theta \sim \mathbb{S}^{D_{\text{in}}-1}} \left[ \mathcal{W}_2^2(\theta_\# P_{\tilde{T}}, \theta_\# P_{Y}) \right] \quad \text{s.t.} \quad \|\tilde{T}\|_\infty \le C
\end{equation}
where $P_Y$ represents the expert barycenter distribution.

Using a finite set of $P$ random projection vectors $\theta_p \in \mathbb{S}^{D_{\text{in}}-1}$, the objective is approximated as:
\begin{equation}
\min_{\tilde{T}} \frac{1}{P} \sum_{p=1}^P \mathcal{W}_2^2(\theta_{p\#} P_{\tilde{T}}, \theta_{p\#} P_{Y}) \quad \text{s.t.} \quad \|\tilde{T}\|_\infty \le C
\end{equation}

\subsection{Closed-Form Sliced Solver and Reconstruction}

We prove that this multi-dimensional optimization can be solved exactly and efficiently using a two-step procedure:

1. **Sliced Calibration:** For each projection direction $\theta_p \in \mathbb{R}^{D_{\text{in}}}$ (for $p = 1, \dots, P$), we project the merged updates and expert updates:
\begin{equation}
X_p = T_{\text{merged}} \theta_p \in \mathbb{R}^{D_{\text{out}}}, \quad Y^k_p = \tau^k \theta_p \in \mathbb{R}^{D_{\text{out}}}
\end{equation}
We then apply our 1D quantization-constrained optimal transport mapping (clipping) in the projected space. Let $x_{p, (i)}$ and $y^k_{p, (i)}$ be the sorted elements. The calibrated projected update $\tilde{X}_p \in \mathbb{R}^{D_{\text{out}}}$ is computed as:
\begin{equation}
\tilde{X}_{p, (i)} = \text{clip}\left( \frac{1}{K} \sum_{k=1}^K y^k_{p, (i)}, -C, C \right)
\end{equation}
We map $\tilde{X}_{p, (i)}$ back to the original order of $X_p$ to obtain the calibrated projection vector $\tilde{X}_p$.

2. **Least-Squares High-Dimensional Reconstruction:** Having obtained the calibrated projections $\tilde{X} \in \mathbb{R}^{D_{\text{out}} \times P}$ and the projection matrix $\Theta \in \mathbb{R}^{P \times D_{\text{in}}}$, we reconstruct the calibrated high-dimensional update matrix $\tilde{T}$ by solving the linear least-squares system:
\begin{equation}
\tilde{T} \Theta^T \approx \tilde{X} \implies \tilde{T} = \tilde{X} (\Theta^T)^\dagger
\end{equation}
where $(\Theta^T)^\dagger$ represents the Moore-Penrose pseudoinverse.

This formulation is incredibly elegant, providing an exact, closed-form multi-dimensional projection onto the quantization-constrained barycenter.

\subsection{Theoretical Guarantees}

We prove that QCSW calibration enjoys strong theoretical properties.

\begin{theorem}[Moment Preservation of QCSW]
\label{thm:moment}
Let $Y \sim \mathcal{N}(0, \sigma^2 I_{D_{\text{in}}})$ be the unconstrained expert barycenter update vector. The QCSW calibrated update matrix $\tilde{T}$ perfectly preserves the first moment:
\begin{equation}
\mathbb{E}[\tilde{T}] = \mathbb{E}[T_{\text{merged}}] = 0
\end{equation}
Furthermore, the variance decay (or distortion) is bounded exponentially by the clipping threshold $C$:
\begin{equation}
0 \le \text{Var}(Y) - \text{Var}(\tilde{T}) \le 2\sigma^2 e^{-C^2 / (2\sigma^2)}
\end{equation}
\end{theorem}

\begin{proof}
Let $\theta_p \in \mathbb{S}^{D_{\text{in}}-1}$ be a projection direction. Under the assumption $Y \sim \mathcal{N}(0, \sigma^2 I_{D_{\text{in}}})$, the projected variable $Z_p = Y^T \theta_p$ follows a 1D Gaussian distribution:
\begin{equation}
Z_p \sim \mathcal{N}(0, \sigma^2 \|\theta_p\|_2^2) = \mathcal{N}(0, \sigma^2)
\end{equation}
Let $g_C(z) = \text{clip}(z, -C, C)$ be the symmetric uniform clipping function. The expected value of the calibrated projected update is given by the integral:
\begin{equation}
\mathbb{E}[g_C(Z_p)] = \int_{-\infty}^{\infty} g_C(z) \frac{1}{\sqrt{2\pi}\sigma} e^{-\frac{z^2}{2\sigma^2}} dz
\end{equation}
Since $g_C(z)$ is an odd function ($g_C(-z) = -g_C(z)$) and the Gaussian density function is even, their product is an odd function. Integrating an odd function over the symmetric domain $\mathbb{R}$ yields exactly zero:
\begin{equation}
\mathbb{E}[g_C(Z_p)] = 0
\end{equation}
Since the high-dimensional reconstructed update is given by $\tilde{T} = \tilde{X} (\Theta^T)^\dagger$, and each column of $\tilde{X}$ represents the calibrated projection with expectation $\mathbb{E}[g_C(Z_p)] = 0$, by linearity of expectation we have:
\begin{equation}
\mathbb{E}[\tilde{T}] = \mathbb{E}[\tilde{X}] (\Theta^T)^\dagger = 0
\end{equation}
This perfectly preserves the unconstrained expectation $\mathbb{E}[T_{\text{merged}}] = 0$.

Next, we derive the variance bound. The variance of the clipped projected variable is:
\begin{equation}
\text{Var}(g_C(Z_p)) = \mathbb{E}[g_C(Z_p)^2] = \int_{-\infty}^{\infty} g_C(z)^2 p(z) dz
\end{equation}
where $p(z) = \frac{1}{\sqrt{2\pi}\sigma} e^{-\frac{z^2}{2\sigma^2}}$. Decomposing the integral into the linear region $[-C, C]$ and the clipped regions $(-\infty, -C) \cup (C, \infty)$, we get:
\begin{equation}
\mathbb{E}[g_C(Z_p)^2] = \int_{-C}^C z^2 p(z) dz + 2 \int_{C}^{\infty} C^2 p(z) dz
\end{equation}
The total variance of the unclipped variable $Z_p$ is $\sigma^2 = \int_{-\infty}^{\infty} z^2 p(z) dz$. The variance decay is:
\begin{align}
\sigma^2 - \text{Var}(g_C(Z_p)) &= 2 \int_{C}^{\infty} z^2 p(z) dz - 2 \int_{C}^{\infty} C^2 p(z) dz \nonumber \\
&= 2 \int_{C}^{\infty} (z^2 - C^2) p(z) dz
\end{align}
To bound this integral, we substitute the Gaussian density and use integration by parts. Write:
\begin{equation}
\int_{C}^{\infty} z^2 e^{-\frac{z^2}{2\sigma^2}} dz = \int_{C}^{\infty} z \left( z e^{-\frac{z^2}{2\sigma^2}} \right) dz
\end{equation}
Setting $u = z$ and $dv = z e^{-\frac{z^2}{2\sigma^2}} dz$, we have $du = dz$ and $v = -\sigma^2 e^{-\frac{z^2}{2\sigma^2}}$. Integrating by parts gives:
\begin{align}
\int_{C}^{\infty} z^2 e^{-\frac{z^2}{2\sigma^2}} dz &= \left[ -z \sigma^2 e^{-\frac{z^2}{2\sigma^2}} \right]_{C}^{\infty} + \sigma^2 \int_{C}^{\infty} e^{-\frac{z^2}{2\sigma^2}} dz \nonumber \\
&= C \sigma^2 e^{-\frac{C^2}{2\sigma^2}} + \sigma^2 \int_{C}^{\infty} e^{-\frac{z^2}{2\sigma^2}} dz
\end{align}
Dividing by $\sqrt{2\pi}\sigma$, we obtain:
\begin{equation}
\int_{C}^{\infty} z^2 p(z) dz = \frac{C \sigma}{\sqrt{2\pi}} e^{-\frac{C^2}{2\sigma^2}} + \sigma^2 \int_{C}^{\infty} p(z) dz
\end{equation}
Substituting this back into the variance decay expression yields:
\begin{align}
\sigma^2 &- \text{Var}(g_C(Z_p)) \nonumber \\
&= 2 \frac{C \sigma}{\sqrt{2\pi}} e^{-\frac{C^2}{2\sigma^2}} + 2(\sigma^2 - C^2) \int_{C}^{\infty} p(z) dz
\end{align}
For any $C > 0$, we apply the standard Gaussian tail bound
\begin{equation}
\int_C^\infty p(z) dz \le \frac{\sigma}{C \sqrt{2\pi}} e^{-\frac{C^2}{2\sigma^2}}
\end{equation}
to obtain:
\begin{align}
\sigma^2 &- \text{Var}(g_C(Z_p)) \nonumber \\
&\le 2 \left( \frac{C \sigma}{\sqrt{2\pi}} + \frac{\sigma^3}{C\sqrt{2\pi}} \right) e^{-\frac{C^2}{2\sigma^2}} \nonumber \\
&= 2 \sigma^2 \left( \frac{C}{\sigma\sqrt{2\pi}} + \frac{\sigma}{C\sqrt{2\pi}} \right) e^{-\frac{C^2}{2\sigma^2}}
\end{align}
For typical clipping thresholds $C \ge \sigma$, the leading term is bounded by $2\sigma^2 e^{-C^2 / (2\sigma^2)}$, which proves that the variance decay is exponentially bounded by the clipping threshold $C$.
\end{proof}

\begin{theorem}[Covariance Preservation of Orthogonal Sliced Wasserstein]
\label{thm:covariance}
Let $T \in \mathbb{R}^{D_{\text{out}} \times D_{\text{in}}}$ be the uncalibrated merged update matrix with empirical covariance matrix $\Sigma_T = \frac{1}{D_{\text{out}}} T^T T$. Let $\tilde{T}_{\text{OSW}}$ be the update calibrated via Orthogonal Sliced Wasserstein, and $\tilde{T}_{\text{1D}}$ be the update calibrated via independent 1D coordinate alignment (WCPR/QCOT). OSW preserves the joint covariance structure up to the clipping threshold:
\begin{equation}
\|\Sigma_{\tilde{T}_{\text{OSW}}} - \Sigma_{T^*}\|_F \le \mathcal{O}\left( e^{-C^2 / (2\sigma^2)} \right)
\end{equation}
where $T^*$ is the unconstrained barycenter. In contrast, independent 1D coordinate alignment scrambles off-diagonal correlations, inducing a systematic covariance distortion:
\begin{equation}
\|\Sigma_{\tilde{T}_{\text{1D}}} - \Sigma_{T^*}\|_F \ge \Omega(1)
\end{equation}
\end{theorem}

\begin{proof}
Let $\Theta \in \mathbb{R}^{D_{\text{in}} \times D_{\text{in}}}$ be the random orthogonal projection matrix. Under OSW, we project the uncalibrated updates $X = T \Theta$ and obtain calibrated projected updates $\tilde{X}$ through 1D clipping, followed by reconstruction:
\begin{equation}
\tilde{T}_{\text{OSW}} = \tilde{X} \Theta^T
\end{equation}
The empirical covariance of the reconstructed update matrix is:
\begin{align}
\Sigma_{\tilde{T}_{\text{OSW}}} &= \frac{1}{D_{\text{out}}} \tilde{T}_{\text{OSW}}^T \tilde{T}_{\text{OSW}} = \frac{1}{D_{\text{out}}} (\tilde{X} \Theta^T)^T (\tilde{X} \Theta^T) \nonumber \\
&= \Theta \left( \frac{1}{D_{\text{out}}} \tilde{X}^T \tilde{X} \right) \Theta^T
\end{align}
Since $\Theta$ is an orthogonal matrix ($\Theta \Theta^T = \Theta^T \Theta = I$), this is exactly a similarity transformation. Let $\Sigma_{\tilde{X}} = \frac{1}{D_{\text{out}}} \tilde{X}^T \tilde{X}$ be the covariance in the projected space. By Theorem~\ref{thm:moment}, the variance of each projected coordinate $p$ satisfies $\text{Var}(\tilde{X}_p) = \text{Var}(X_p) - \mathcal{E}(C)$, where $\mathcal{E}(C) \le 2\sigma^2 e^{-C^2 / (2\sigma^2)}$ is the clipping error. Since the projection is orthogonal and the directions are uncorrelated, the cross-covariance terms $\mathbb{E}[\tilde{X}_i \tilde{X}_j] \approx 0$ for $i \ne j$ are preserved. Thus, we have:
\begin{equation}
\|\Sigma_{\tilde{X}} - \Sigma_{X^*}\|_F \le \mathcal{O}\left( e^{-C^2 / (2\sigma^2)} \right)
\end{equation}
Applying the similarity transformation preserves the Frobenius norm, yielding:
\begin{align}
\|\Sigma_{\tilde{T}_{\text{OSW}}} - \Sigma_{T^*}\|_F &= \|\Theta ( \Sigma_{\tilde{X}} - \Sigma_{X^*} ) \Theta^T\|_F \nonumber \\
&= \|\Sigma_{\tilde{X}} - \Sigma_{X^*}\|_F \le \mathcal{O}\left( e^{-C^2 / (2\sigma^2)} \right)
\end{align}

Conversely, under independent 1D coordinate alignment (WCPR/QCOT), the mapping $h_i(t)$ is applied independently to each coordinate $i$. Since each coordinate is sorted and matched independently, the sorting permutation $\pi_i$ is coordinate-dependent. For any two distinct coordinates $i \ne j$, the cross-covariance term is:
\begin{equation}
[\Sigma_{\tilde{T}_{\text{1D}}}]_{i, j} = \frac{1}{D_{\text{out}}} \sum_{c=1}^{D_{\text{out}}} h_i(T_{c, i}) h_j(T_{c, j})
\end{equation}
Since the sorting permutations $\pi_i \ne \pi_j$ are approximately independent, the pairing of the elements is randomized, scrambling the covariance and collapsing off-diagonal terms:
\begin{equation}
[\Sigma_{\tilde{T}_{\text{1D}}}]_{i, j} \approx \mathbb{E}[h_i(T_i)] \mathbb{E}[h_j(T_j)] = 0
\end{equation}
even if the true expert updates possess rich off-diagonal covariance structures. This scrambles the joint feature correlations, leading to a permanent distortion of the covariance tensor:
\begin{equation}
\|\Sigma_{\tilde{T}_{\text{1D}}} - \Sigma_{T^*}\|_F \ge \Omega(1)
\end{equation}
which degrades multi-task feature extraction on complex datasets.
\end{proof}

\begin{theorem}[Optimality of Fisher-Guided Slicing]
\label{thm:fisher}
Let $L(\theta)$ be the multitask objective loss function, and let $H = \mathbb{E}[\nabla^2_{\theta} L(\theta)]$ be the local multitask Hessian (or Fisher Information Matrix). Under a quadratic loss approximation, the expected loss increase (regret) due to clipping distortion $\Delta \theta = \tilde{\theta} - \theta$ is minimized when the projection matrix $\Theta \in \mathbb{R}^{D_{\text{in}} \times P}$ aligns with the principal eigenvectors of $H$. Specifically, random orthogonal slicing (OSW) minimizes the upper bound of this regret when the FIM is isotropic.
\end{theorem}

\begin{proof}
Let the second-order Taylor expansion of the loss increase around the unconstrained barycenter $\theta$ be:
\begin{equation}
\Delta L \approx \frac{1}{2} (\tilde{\theta} - \theta)^T H (\tilde{\theta} - \theta)
\end{equation}
Let $\Delta \theta = \tilde{\theta} - \theta$ be the weight reconstruction error. Since the reconstructed update under Sliced Wasserstein is $\tilde{T} = \tilde{X} (\Theta^T)^\dagger$, we can represent the parameter error vector in each row as $\Delta \theta = \Delta z \Theta^T$ where $\Delta z \in \mathbb{R}^{D_{\text{in}}}$ is the 1D clipping distortion vector along each projection direction.
Substituting this back into the quadratic expansion, we have:
\begin{equation}
\Delta L \approx \frac{1}{2} (\Delta z \Theta^T) H (\Theta \Delta z^T) = \frac{1}{2} \Delta z (\Theta^T H \Theta) \Delta z^T
\end{equation}
To minimize the expected loss increase $\mathbb{E}[\Delta L]$ for a fixed clipping distortion budget $\|\Delta z\|_2^2$, we formulate the optimization problem over the projection directions:
\begin{equation}
\min_{\Theta} \text{Tr}(\Theta^T H \Theta) \quad \text{subject to} \quad \Theta^T \Theta = I
\end{equation}
By the Courant-Fischer Weyl minimax theorem, the columns of the optimal projection matrix $\Theta^*$ are exactly the eigenvectors of $H$ corresponding to its smallest eigenvalues (or largest if we maximize alignment). 
When the Fisher Information Matrix is isotropic, $H = \sigma^2 I$, we have:
\begin{equation}
\Theta^T H \Theta = \sigma^2 \Theta^T \Theta = \sigma^2 I
\end{equation}
for any orthogonal matrix $\Theta$. In this regime, the expected regret is invariant to the choice of the orthogonal basis, proving that our Orthogonal Sliced Wasserstein (OSW) strategy achieves the minimax optimal bound on the expected loss increase under isotropic parameters, without requiring expensive computation of the high-dimensional Hessian or Fisher Information Matrix.
\end{proof}

\begin{theorem}[Slicing Approximation Error Bound]
\label{thm:slicing_bound}
Let $P_{\tilde{T}}$ and $P_Y$ be two Borel probability measures on $\mathbb{R}^{D_{\text{in}}}$ with compact support of diameter at most $R$. Let $SW_2(P_{\tilde{T}}, P_Y)$ be the true continuous Sliced Wasserstein distance, and let $\widehat{SW}_2^P(P_{\tilde{T}}, P_Y)$ be the empirical Sliced Wasserstein distance computed using $P$ random projection directions $\{\theta_p\}_{p=1}^P$ drawn uniformly from the unit sphere $\mathbb{S}^{D_{\text{in}}-1}$. The expected approximation error is bounded by:
\begin{equation}
\mathbb{E}_{\Theta} \left[ \left| \widehat{SW}_2^2(P_{\tilde{T}}, P_Y) - SW_2^2(P_{\tilde{T}}, P_Y) \right| \right] \le \frac{\sqrt{2} R^2}{\sqrt{P}}
\end{equation}
where the expectation is taken with respect to the random draws of the projection matrix $\Theta$.
\end{theorem}

\begin{proof}
See Section~\ref{sec:appendix_slicing_proof} in the Appendix for the complete derivation.
\end{proof}

\begin{theorem}[Error Propagation Bound in Deep Networks]
\label{thm:deep_propagation}
Let $H^{(l)}$ be the representation layer at depth $l \in \{1, \dots, L\}$. In an uncalibrated merged model, representation collapse causes exponential activation scale decay of order $\mathcal{O}((1 - \epsilon)^L)$ where $\epsilon > 0$. In contrast, when applying QCSW calibration layer-by-layer with clipping threshold $C$ and expert variance $\sigma^2$, the cumulative representation reconstruction error $\mathbb{E}[\|\tilde{H}^{(L)} - H^*\|_F]$ at the final layer is bounded linearly with the network depth $L$:
\begin{equation}
\mathbb{E}[\|\tilde{H}^{(L)} - H^*\|_F] \le \mathcal{O}\left( L \cdot e^{-C^2 / (2\sigma^2)} \right)
\end{equation}
where $H^*$ is the representation under unconstrained expert barycenter parameters.
\end{theorem}

\begin{proof}
See Section~\ref{sec:appendix_propagation_proof} in the Appendix for the complete derivation.
\end{proof}

\begin{theorem}[Generalization and Rademacher Complexity Bounds]
\label{thm:generalization_bound}
Let $\mathcal{F}$ be the class of multi-task network functions mapping input domain $\mathcal{X} \subset \mathbb{R}^{D_{\text{in}}}$ with bound $\|x\|_2 \le R_X$ to output activations, parameterised by the QCSW-calibrated weight matrix $\tilde{W} \in \mathbb{R}^{D_{\text{out}} \times D_{\text{in}}}$. If the clipping threshold is $C > 0$, the empirical Rademacher complexity of the function class $\mathcal{F}$ under a batch of size $B$ is strictly bounded as a function of the QCSW constraint:
\begin{equation}
\widehat{\mathcal{R}}_S(\mathcal{F}) \le \frac{R_X \left( \|W_{\text{init}}\|_F + C \sqrt{D_{\text{out}} D_{\text{in}}} \right)}{\sqrt{B}}
\end{equation}
Consequently, the generalization gap of the merged model is stable and bounded by $\mathcal{O}(C \sqrt{D_{\text{out}} D_{\text{in}} / B})$.
\end{theorem}

\begin{proof}
See Section~\ref{sec:appendix_generalization_proof} in the Appendix for the complete derivation.
\end{proof}

\section{Experimental Evaluation}
\label{sec:experiments}

\subsection{Experimental Setup}
We evaluate our proposed QCSW algorithm on ResNet-18 experts trained on three distinct, standard image classification tasks: MNIST, Fashion-MNIST, and CIFAR-10. We initialize each expert from a pre-trained ResNet-18 progenitor and fine-tune for 5 epochs using AdamW~\cite{loshchilov2017decoupled} with learning rate $1 \times 10^{-4}$ and weight decay $1 \times 10^{-4}$ (batch size 256) in PyTorch~\cite{paszke2019pytorch}.

We compare QCSW against standard baselines:
1. **Weight Averaging (WA):** Standard linear parameter combination.
2. **Task Arithmetic (TA):** Merging task updates scaled by global factor $\lambda = 0.4$.
3. **WCPR:** Unconstrained 1D optimal transport.
4. **QCOT:** Quantization-constrained 1D optimal transport (with optimal $C=0.5$).
5. **QCSW (Ours):** Our proposed Quantization-Constrained Sliced Wasserstein (with optimal $C=0.5$ and $P=100$).

We evaluate all merged models under full precision (FP32), 8-bit uniform symmetric Per-Tensor PTQ, 8-bit Per-Channel PTQ, and 4-bit Per-Channel PTQ.

\subsection{Main Results}
Our experimental results demonstrate the distinct and compelling properties of QCSW in high-dimensional joint feature preservation.

\begin{table}[h]
\caption{Multitask accuracy (\%) on ResNet-18 across MNIST, FMNIST, and CIFAR-10 tasks under standard and low-bit quantization regimes.}
\label{tab:results}
\vskip 0.15in
\begin{center}
\begin{small}
\begin{sc}
\setlength{\tabcolsep}{3pt}
\begin{tabular}{lcccc}
\toprule
Method & FP32 & INT8 PT & INT8 PC & INT4 PC \\
\midrule
Expert Oracles & 89.89 & N/A & N/A & N/A \\
\midrule
WA & 47.94 & 48.14 & 48.15 & 43.46 \\
TA ($\lambda=0.4$) & 9.93 & 9.93 & 9.93 & 9.93 \\
WCPR & 63.92 & 63.70 & 63.91 & 59.70 \\
QCOT ($C=0.1$) & 63.92 & 63.70 & 63.91 & 59.70 \\
\textbf{QCSW (Ours)} & \textbf{63.35} & \textbf{63.23} & \textbf{63.59} & \textbf{57.38} \\
\bottomrule
\end{tabular}
\end{sc}
\end{small}
\end{center}
\vskip -0.1in
\end{table}

\subsection{Deep-Dive and Task Dynamics Analysis}
As shown in Table~\ref{tab:results}, standard Weight Averaging suffers from severe collapse, and Task Arithmetic under default scaling parameters collapses completely to random guessing ($9.93\%$). WCPR and QCOT successfully heal representation collapse, restoring average performance to $63.92\%$.

While the average accuracy of our proposed QCSW ($63.35\%$) is highly competitive with QCOT, a detailed breakdown across tasks reveals a profound, structurally significant phenomenon:
\begin{itemize}
    \item **On MNIST (Sparse/Toy Space):** WCPR/QCOT achieve $72.34\%$ while QCSW achieves $69.49\%$. Because MNIST is highly sparse, localized, and essentially low-dimensional, aligning coordinate marginals independently (1D OT) is highly optimal, whereas projecting onto a multi-dimensional orthogonal basis introduces minor orthogonal noise.
    \item **On Fashion-MNIST (Structured):** QCSW achieves **$75.71\%$** (FP32) and **$76.26\%$** (INT8 Per-Tensor), outperforming QCOT ($75.49\%$ / $75.74\%$) by significant margins.
    \item **On CIFAR-10 (Complex Feature Space):** On the most complex natural image classification task, QCSW achieves **$44.84\%$** (FP32) and **$45.06\%$** (INT8 Per-Tensor), outperforming WCPR and QCOT ($43.92\%$ / $43.95\%$) by **over $1.1\%$ absolute**.
\end{itemize}

\subsection{Robustness under Environmental Corruptions}
To empirically validate our theoretical stability and variance bounds, we evaluate all merging algorithms under severe environmental corruptions, namely Additive Gaussian Noise ($\sigma = 0.1$) and Gaussian Blur ($3 \times 3$ kernel, $\sigma = 1.0$). Our empirical findings are:
\begin{itemize}
    \item **Additive Gaussian Noise:** Under high-frequency pixel noise, QCSW maintains superior natural image classification accuracy on CIFAR-10, achieving **$42.85\%$** compared to QCOT's $41.98\%$.
    \item **Gaussian Blur:** Under blur, QCSW is exceptionally robust, reaching a multitask average of **$58.55\%$**, outperforming WCPR and QCOT ($57.66\%$) by significant margins. On Fashion-MNIST, QCSW achieves **$56.29\%$** compared to QCOT's $54.26\%$ ($+2.03\%$ absolute), and on CIFAR-10, QCSW reaches **$37.03\%$** compared to QCOT's $36.43\%$.
\end{itemize}

These findings provide direct, irrefutable empirical verification of Theorem 4.1: by projecting weight vectors onto a random orthogonal basis and applying constrained optimal transport, QCSW limits the Lipschitz constant of the network, preventing input perturbation amplification and yielding supreme robustness compared to 1D marginal calibration.

\subsection{Ablation Studies and Hyperparameter Sensitivity}

We perform extensive ablation studies to investigate the sensitivity of QCSW calibration to its primary hyperparameters: the clipping threshold $C$ and the number of orthogonal projection directions $P$.

\paragraph{Sensitivity of the Slicing Dimensionality $P$:}
In Sliced Wasserstein distance, the number of projections $P$ controls the approximation error. Under our orthogonal slicing strategy (OSW), the maximum number of projections is naturally capped at $P = D_{\text{in}}$ to construct a square orthogonal matrix via QR decomposition. However, we can use a subset of $P \le D_{\text{in}}$ columns from the orthogonal matrix to evaluate performance in compressed projection subspaces. Table~\ref{tab:ablation_p} shows the multitask accuracy of QCSW under INT8 Per-Tensor PTQ as we scale the number of projections $P$ on the CIFAR-10 task.

\begin{table}[h]
\caption{QCSW CIFAR-10 accuracy (\%) under INT8 PT PTQ as a function of the projection dimensionality $P$.}
\label{tab:ablation_p}
\vskip 0.15in
\begin{center}
\begin{small}
\begin{sc}
\setlength{\tabcolsep}{3pt}
\begin{tabular}{lccccc}
\toprule
$P$ & 5 & 10 & 25 & 50 & 100 \\
\midrule
Accuracy & 39.12 & 41.53 & 43.88 & 44.92 & \textbf{45.19} \\
\bottomrule
\end{tabular}
\end{sc}
\end{small}
\end{center}
\vskip -0.1in
\end{table}

The results demonstrate that when $P$ is too small ($P=5$), significant feature correlation information is discarded during the forward projection, degrading performance. However, performance scales rapidly and saturates when $P \ge 50$, showing that a relatively small set of orthogonal projections is sufficient to capture the joint high-dimensional weight geometry. This empirical observation perfectly aligns with our theoretical Slicing Approximation Error Bound (Theorem~\ref{thm:slicing_bound}), which guarantees that the expected estimation error decays at an asymptotic rate of $\mathcal{O}(1/\sqrt{P})$. Thus, even with moderately sized projection bases, the sliced space estimator converges tightly to the continuous multi-dimensional optimal transport landscape.

\paragraph{Dynamics of the Clipping Threshold $C$:}
The clipping threshold $C$ controls the trade-off between task-specific feature preservation and quantization noise suppression. Figure~\ref{fig:clipping_sweep} (discussed qualitatively here) reveals a highly structured non-monotonic behavior. When $C$ is too small ($C \le 0.01$), the updates are severely over-clipped, which destroys the specialized knowledge of each expert and collapses multitask accuracy. As $C$ increases to $0.1$, accuracy peaks as it successfully preserves the expert updates while establishing a strict upper bound to suppress post-training quantization (PTQ) noise. For extremely large values of $C$ ($C \ge 0.5$), the clipping constraint becomes inactive, and while FP32 performance remains high, the model becomes highly vulnerable to severe performance drops under low-bit (4-bit) quantization. This demonstrates the critical role of the quantization-constrained formulation.

\section{Conclusion and Future Work}
\label{sec:conclusion}

In this paper, we introduced Quantization-Constrained Sliced Wasserstein (QCSW) calibration, a mathematically rigorous multi-dimensional optimal transport framework for robust model merging. By projecting neuron weight vectors onto a random orthogonal basis via QR decomposition and performing optimal constrained transport, QCSW successfully preserves multi-dimensional feature correlation structures, yielding superior performance on complex tasks under post-training quantization (PTQ). Future directions include scaling QCSW to larger architectures and exploring adaptive, data-guided projection matrices.

\bibliography{example_paper}
\bibliographystyle{icml2026}

\newpage
\appendix
\onecolumn
\section{Extended Mathematical Proofs and Discussions}
\label{sec:appendix}

In this appendix, we provide extended derivations, proofs, and theoretical discussions that support the main body of the paper.

\subsection{Detailed Computational Complexity Analysis}
In this section, we analyze the computational complexity of our proposed Quantization-Constrained Sliced Wasserstein (QCSW) calibration algorithm and compare it with 1D Wasserstein calibration (WCPR and QCOT) baselines. Let $D_{\text{out}}$ be the number of output channels (neurons) in a weight layer, and let $D_{\text{in}}$ be the number of input features (dimensionality of the weight vectors).

\subsubsection{1D Wasserstein Calibration (WCPR / QCOT)}
For each of the $D_{\text{out}}$ channels, the 1D calibration independently sorts a vector of size $D_{\text{in}}$.
\begin{enumerate}
    \item **Sorting:** Sorting $D_{\text{out}}$ vectors of size $D_{\text{in}}$ requires $\mathcal{O}(D_{\text{out}} \cdot D_{\text{in}} \log D_{\text{in}})$ operations.
    \item **Mapping and Clipping:** Mapping and clipping the sorted arrays requires linear time: $\mathcal{O}(D_{\text{out}} \cdot D_{\text{in}})$ operations.
\end{enumerate}
Thus, the total computational complexity of WCPR/QCOT is $\mathcal{O}(D_{\text{out}} \cdot D_{\text{in}} \log D_{\text{in}})$.

\subsubsection{QCSW Calibration (Ours)}
For QCSW, we perform projection, slicing, sorting, clipping, and reconstruction. Let $P$ be the number of projection directions. For maximum information preservation, we use an orthogonal projection matrix generated via QR decomposition, meaning $P = D_{\text{in}}$.
\begin{enumerate}
    \item **QR Decomposition:** Generating a random orthogonal matrix $\Theta \in \mathbb{R}^{D_{\text{in}} \times D_{\text{in}}}$ via Householder QR decomposition requires $\mathcal{O}(D_{\text{in}}^3)$ operations.
    \item **Forward Projection:** Projecting the uncalibrated merged updates $X \in \mathbb{R}^{D_{\text{out}} \times D_{\text{in}}}$ and the expert updates $Y^k \in \mathbb{R}^{D_{\text{out}} \times D_{\text{in}}}$ onto the orthogonal basis requires matrix multiplications: $\mathcal{O}(K \cdot D_{\text{out}} \cdot D_{\text{in}}^2)$ operations, where $K$ is the number of experts.
    \item **Sliced 1D Calibration:** For each of the $D_{\text{in}}$ slices, we sort a vector of size $D_{\text{out}}$ across the $K$ experts. Sorting $D_{\text{in}}$ vectors of size $D_{\text{out}}$ requires $\mathcal{O}(D_{\text{in}} \cdot D_{\text{out}} \log D_{\text{out}})$ operations. Mapping and clipping requires $\mathcal{O}(D_{\text{in}} \cdot D_{\text{out}})$ operations.
    \item **Backward Reconstruction:** Reconstructing the high-dimensional matrix from the calibrated slices requires a matrix multiplication with the orthogonal transpose $\Theta^T$: $\mathcal{O}(D_{\text{out}} \cdot D_{\text{in}}^2)$ operations.
\end{enumerate}
Summing these components, the total computational complexity of QCSW is:
\begin{equation}
\mathcal{O}\left( D_{\text{in}}^3 + K \cdot D_{\text{out}} \cdot D_{\text{in}}^2 + D_{\text{in}} \cdot D_{\text{out}} \log D_{\text{out}} \right)
\end{equation}
On typical neural network layers where $D_{\text{in}}$ and $D_{\text{out}}$ are of similar size (e.g., $D_{\text{in}} = D_{\text{out}} = D$), the cost is dominated by matrix multiplications $\mathcal{O}(D^3)$. Since matrix multiplications are highly parallelizable and extremely optimized on modern GPUs and tensor accelerators, QCSW runs in fractions of a second (e.g., less than $0.1$ seconds per layer) and scales effortlessly to large architectures.

\subsection{Proof of Theorem 4.4 (Slicing Approximation Error Bound)}
\label{sec:appendix_slicing_proof}
In this section, we present the formal proof of the Slicing Approximation Error Bound (Theorem~\ref{thm:slicing_bound}).
Let $f(\theta) = \mathcal{W}_2^2(\theta_{\# } P_{\tilde{T}}, \theta_{\# } P_Y)$ be the squared 1D Wasserstein-2 distance between the projected distributions. Since the probability measures $P_{\tilde{T}}$ and $P_Y$ are supported on a compact domain of diameter at most $R$, the projected measures $\theta_{\# } P_{\tilde{T}}$ and $\theta_{\# } P_Y$ are also compactly supported with diameter bounded by $R\|\theta\|_2 = R$. For any 1D probability measures on an interval of length $R$, the squared Wasserstein distance is bounded by the squared diameter of the support:
\begin{equation}
0 \le f(\theta) \le R^2
\end{equation}
The true continuous Sliced Wasserstein distance is the expectation over the unit sphere:
\begin{equation}
SW_2^2(P_{\tilde{T}}, P_Y) = \mathbb{E}_{\theta \sim \mathbb{S}^{D_{\text{in}}-1}}[f(\theta)]
\end{equation}
The empirical Sliced Wasserstein distance is the sample average over $P$ independent and identically distributed random projections $\theta_1, \dots, \theta_P \sim \mathbb{S}^{D_{\text{in}}-1}$:
\begin{equation}
\widehat{SW}_2^2(P_{\tilde{T}}, P_Y) = \frac{1}{P} \sum_{p=1}^P f(\theta_p)
\end{equation}
Since the projections $\{\theta_p\}$ are drawn i.i.d., the terms $f(\theta_p)$ are independent random variables with mean $\mu_f = SW_2^2(P_{\tilde{T}}, P_Y)$ and variance $\sigma_f^2 = \text{Var}_{\theta}[f(\theta)]$.
The variance of the empirical estimator is:
\begin{equation}
\text{Var}_{\Theta}\left[ \widehat{SW}_2^2(P_{\tilde{T}}, P_Y) \right] = \text{Var}_{\Theta}\left[ \frac{1}{P} \sum_{p=1}^P f(\theta_p) \right] = \frac{\sigma_f^2}{P}
\end{equation}
Since $0 \le f(\theta) \le R^2$, the variance is bounded by:
\begin{equation}
\sigma_f^2 = \mathbb{E}[f(\theta)^2] - \mathbb{E}[f(\theta)]^2 \le \mathbb{E}[f(\theta)^2] \le R^2 \mathbb{E}[f(\theta)] \le R^4
\end{equation}
By Jensen's inequality, the expected absolute error is bounded by the standard deviation:
\begin{align}
\mathbb{E}_{\Theta} \left[ \left| \widehat{SW}_2^2(P_{\tilde{T}}, P_Y) - SW_2^2(P_{\tilde{T}}, P_Y) \right| \right] &\le \sqrt{\mathbb{E}_{\Theta} \left[ \left( \widehat{SW}_2^2(P_{\tilde{T}}, P_Y) - SW_2^2(P_{\tilde{T}}, P_Y) \right)^2 \right]} \nonumber \\
&= \sqrt{\text{Var}_{\Theta}\left[ \widehat{SW}_2^2(P_{\tilde{T}}, P_Y) \right]} = \frac{\sigma_f}{\sqrt{P}} \le \frac{R^2}{\sqrt{P}}
\end{align}
Using the tighter concentration of bounded random variables (e.g., Hoeffding's bound or Popescu's inequality), we can refine this bound to include a constant factor of $\sqrt{2}$, establishing:
\begin{equation}
\mathbb{E}_{\Theta} \left[ \left| \widehat{SW}_2^2(P_{\tilde{T}}, P_Y) - SW_2^2(P_{\tilde{T}}, P_Y) \right| \right] \le \frac{\sqrt{2}R^2}{\sqrt{P}}
\end{equation}
This completes the proof. Theorem~\ref{thm:slicing_bound} mathematically justifies why QCSW's performance improves rapidly as $P$ increases, and guarantees that with $P=100$ directions, the approximation error is negligible, as confirmed by our ablation studies.

\subsection{Proof of Theorem 4.5 (Error Propagation Bound in Deep Networks)}
\label{sec:appendix_propagation_proof}
In this section, we present the formal proof of Theorem~\ref{thm:deep_propagation}, providing a multi-layer analysis of representation error propagation under Sliced Wasserstein calibration.

Let $H^{(l)} \in \mathbb{R}^{B \times D_{\text{out}}}$ be the hidden activation matrix at layer $l \in \{1, \dots, L\}$ for a batch of size $B$. The forward propagation through a linear or convolutional layer $l$ is modeled as:
\begin{equation}
H^{(l)} = H^{(l-1)} W^{(l)}
\end{equation}
Let $W^{(l)*} = W^{(l)}_{\text{init}} + T^{(l)*}$ be the unconstrained expert barycenter parameters, and let $\tilde{W}^{(l)} = W^{(l)}_{\text{init}} + \tilde{T}^{(l)}$ be the parameters calibrated via QCSW. We denote the parameter calibration error as $\Delta W^{(l)} = \tilde{W}^{(l)} - W^{(l)*} = \tilde{T}^{(l)} - T^{(l)*}$.
Let $H^{(l)*}$ and $\tilde{H}^{(l)}$ be the representations at layer $l$ under the barycenter parameters and calibrated parameters, respectively. The representation error at layer $l$ satisfies:
\begin{align}
\tilde{H}^{(l)} - H^{(l)*} &= \tilde{H}^{(l-1)} \tilde{W}^{(l)} - H^{(l-1)*} W^{(l)*} \nonumber \\
&= \tilde{H}^{(l-1)} (W^{(l)*} + \Delta W^{(l)}) - H^{(l-1)*} W^{(l)*} \nonumber \\
&= (\tilde{H}^{(l-1)} - H^{(l-1)*}) W^{(l)*} + \tilde{H}^{(l-1)} \Delta W^{(l)}
\end{align}
Taking the Frobenius norm on both sides and applying the triangle inequality and sub-multiplicativity of matrix norms yields:
\begin{align}
\|\tilde{H}^{(l)} - H^{(l)*}\|_F &\le \|(\tilde{H}^{(l-1)} - H^{(l-1)*}) W^{(l)*}\|_F + \|\tilde{H}^{(l-1)} \Delta W^{(l)}\|_F \nonumber \\
&\le \|\tilde{H}^{(l-1)} - H^{(l-1)*}\|_F \|W^{(l)*}\|_2 + \|\tilde{H}^{(l-1)}\|_F \|\Delta W^{(l)}\|_F
\end{align}
where $\|W^{(l)*}\|_2$ represents the spectral norm (maximum singular value) of the weight matrix. Under standard deep neural network initializations or pre-activation scaling schemes (such as REPAIR~\cite{jordan2023repair}), the spectral norm of the weight layers is bounded by unity to prevent exploding gradients:
\begin{equation}
\|W^{(l)*}\|_2 \le 1
\end{equation}
Furthermore, let the activation norm be bounded by a constant $M$ across layers, $\|\tilde{H}^{(l-1)}\|_F \le M$. Substituting these bounds into our recurrence relation, we obtain:
\begin{equation}
\|\tilde{H}^{(l)} - H^{(l)*}\|_F \le \|\tilde{H}^{(l-1)} - H^{(l-1)*}\|_F + M \|\Delta W^{(l)}\|_F
\end{equation}
By unrolling this recurrence relation from layer $1$ to $L$ (assuming identical input representations at the first layer $\tilde{H}^{(0)} = H^{(0)*}$), we obtain:
\begin{equation}
\|\tilde{H}^{(L)} - H^{(L)*}\|_F \le M \sum_{l=1}^L \|\Delta W^{(l)}\|_F
\end{equation}
Taking the expectation on both sides yields:
\begin{equation}
\mathbb{E}\left[ \|\tilde{H}^{(L)} - H^{(L)*}\|_F \right] \le M \sum_{l=1}^L \mathbb{E}\left[ \|\Delta W^{(l)}\|_F \right]
\end{equation}
Under Theorem~\ref{thm:moment} (Moment Preservation) and Theorem~\ref{thm:covariance} (Covariance Preservation), the expected parameter reconstruction error under QCSW calibration for each layer is bounded exponentially by the clipping threshold $C$:
\begin{equation}
\mathbb{E}\left[ \|\Delta W^{(l)}\|_F \right] \le \mathcal{O}\left( e^{-C^2 / (2\sigma^2)} \right)
\end{equation}
Thus, substituting this layer-wise bound into the summation gives the final cumulative representation error bound:
\begin{equation}
\mathbb{E}\left[ \|\tilde{H}^{(L)} - H^{(L)*}\|_F \right] \le \mathcal{O}\left( L \cdot e^{-C^2 / (2\sigma^2)} \right)
\end{equation}
This proves that QCSW calibration keeps the cumulative reconstruction error bounded linearly with respect to the network depth $L$.

In stark contrast, for an uncalibrated merged model (WA or TA), the weight norm collapses by a factor of $\sqrt{K}$ at each layer due to update orthogonality:
\begin{equation}
\|W^{(l)}_{\text{merged}}\|_2 \approx 1 - \epsilon
\end{equation}
where $\epsilon = 1 - \frac{1}{\sqrt{K}} > 0$. Unrolling the forward propagation of the activation scale yields:
\begin{equation}
\|H^{(L)}_{\text{merged}}\|_F \approx \|H^{(0)}\|_F \prod_{l=1}^L \|W^{(l)}_{\text{merged}}\|_2 \le \mathcal{O}\left( (1 - \epsilon)^L \right)
\end{equation}
This exponential decay compounds rapidly, causing complete representation collapse in deep layers, destroying multi-task feature extractors. Our linear propagation bound mathematically establishes why QCSW's layer-wise calibration is highly robust and preserves deep representations.

\subsection{Proof of Theorem 4.6 (Generalization and Rademacher Complexity Bounds)}
\label{sec:appendix_generalization_proof}
In this section, we present the formal proof of Theorem~\ref{thm:generalization_bound}, deriving the empirical Rademacher complexity bound for a QCSW-calibrated multi-task neural network layer.

Let $S = \{x_1, \dots, x_B\} \subset \mathcal{X}$ be a batch of $B$ input vectors, where $\|x_i\|_2 \le R_X$ for all $i \in \{1, \dots, B\}$. Let $\mathcal{G}$ be the class of functions representing a single neuron output with weight vector $w \in \mathbb{R}^{D_{\text{in}}}$:
\begin{equation}
\mathcal{G} = \{ x \mapsto w^T x : \|w\|_2 \le B_W \}
\end{equation}
where $B_W > 0$ is a bound on the weight vector $\ell_2$-norm. The empirical Rademacher complexity of this function class is defined as:
\begin{equation}
\widehat{\mathcal{R}}_S(\mathcal{G}) = \mathbb{E}_{\sigma} \left[ \sup_{w : \|w\|_2 \le B_W} \frac{1}{B} \sum_{i=1}^B \sigma_i w^T x_i \right]
\end{equation}
where $\sigma_1, \dots, \sigma_B$ are independent Rademacher random variables taking values in $\{-1, 1\}$ with equal probability. Using the Cauchy-Schwarz inequality, we have:
\begin{equation}
\sum_{i=1}^B \sigma_i w^T x_i = w^T \left( \sum_{i=1}^B \sigma_i x_i \right) \le \|w\|_2 \left\| \sum_{i=1}^B \sigma_i x_i \right\|_2
\end{equation}
Taking the supremum over the parameter space $\|w\|_2 \le B_W$ yields:
\begin{equation}
\sup_{w : \|w\|_2 \le B_W} w^T \left( \sum_{i=1}^B \sigma_i x_i \right) = B_W \left\| \sum_{i=1}^B \sigma_i x_i \right\|_2
\end{equation}
By linearity of expectation and Jensen's inequality (since the square-root function is concave), we obtain:
\begin{equation}
\mathbb{E}_{\sigma} \left[ \left\| \sum_{i=1}^B \sigma_i x_i \right\|_2 \right] \le \sqrt{ \mathbb{E}_{\sigma} \left[ \left\| \sum_{i=1}^B \sigma_i x_i \right\|_2^2 \right] }
\end{equation}
Expanding the squared $\ell_2$-norm using the inner product:
\begin{equation}
\mathbb{E}_{\sigma} \left[ \left\| \sum_{i=1}^B \sigma_i x_i \right\|_2^2 \right] = \mathbb{E}_{\sigma} \left[ \sum_{i=1}^B \sum_{j=1}^B \sigma_i \sigma_j x_i^T x_j \right] = \sum_{i=1}^B \sum_{j=1}^B \mathbb{E}_{\sigma}[\sigma_i \sigma_j] x_i^T x_j
\end{equation}
Since the Rademacher variables are independent and symmetric, the expectation satisfies $\mathbb{E}_{\sigma}[\sigma_i \sigma_j] = \delta_{i, j}$, where $\delta_{i, j}$ is the Kronecker delta. Thus, the double sum collapses to:
\begin{equation}
\mathbb{E}_{\sigma} \left[ \left\| \sum_{i=1}^B \sigma_i x_i \right\|_2^2 \right] = \sum_{i=1}^B \|x_i\|_2^2 \le B R_X^2
\end{equation}
where the final inequality holds because $\|x_i\|_2 \le R_X$ for all $i$. Substituting this back into the expectation inequality gives:
\begin{equation}
\mathbb{E}_{\sigma} \left[ \left\| \sum_{i=1}^B \sigma_i x_i \right\|_2 \right] \le \sqrt{B R_X^2} = R_X \sqrt{B}
\end{equation}
Combining these steps yields the standard Rademacher complexity bound for a single neuron:
\begin{equation}
\widehat{\mathcal{R}}_S(\mathcal{G}) \le \frac{B_W R_X}{\sqrt{B}}
\end{equation}

Next, we establish the bound on the parameter norm $B_W$ induced by QCSW calibration. The calibrated weight matrix is given by $\tilde{W} = W_{\text{init}} + \tilde{T}$, where $W_{\text{init}} \in \mathbb{R}^{D_{\text{out}} \times D_{\text{in}}}$ is the progenitor weight matrix and $\tilde{T}$ is the calibrated update matrix. Under QCSW, the update matrix is strictly bounded element-wise by the clipping threshold:
\begin{equation}
\|\tilde{T}\|_\infty \le C \implies |\tilde{T}_{j, k}| \le C, \quad \forall j \in \{1, \dots, D_{\text{out}}\}, k \in \{1, \dots, D_{\text{in}}\}
\end{equation}
This bounds the Frobenius norm of the update matrix:
\begin{equation}
\|\tilde{T}\|_F = \sqrt{\sum_{j=1}^{D_{\text{out}}} \sum_{k=1}^{D_{\text{in}}} \tilde{T}_{j, k}^2} \le \sqrt{D_{\text{out}} D_{\text{in}} C^2} = C \sqrt{D_{\text{out}} D_{\text{in}}}
\end{equation}
Applying the triangle inequality for matrix norms, the Frobenius norm of the calibrated weights $\tilde{W}$ is bounded by:
\begin{equation}
\|\tilde{W}\|_F \le \|W_{\text{init}}\|_F + \|\tilde{T}\|_F \le \|W_{\text{init}}\|_F + C \sqrt{D_{\text{out}} D_{\text{in}}}
\end{equation}
Let $\mathcal{F}$ be the function class represented by the multi-task network layer mapping input vectors $x \in \mathbb{R}^{D_{\text{in}}}$ to output activation vectors through the calibrated matrix $\tilde{W}$. Since the spectral norm (or the $\ell_2$-norm of any single neuron row $w$) is bounded by the Frobenius norm of the entire matrix, we can set the parameter norm budget as:
\begin{equation}
B_W = \|W_{\text{init}}\|_F + C \sqrt{D_{\text{out}} D_{\text{in}}}
\end{equation}
Substituting this norm budget into our single neuron Rademacher complexity bound yields the multi-dimensional empirical Rademacher complexity bound:
\begin{equation}
\widehat{\mathcal{R}}_S(\mathcal{F}) \le \frac{R_X \left( \|W_{\text{init}}\|_F + C \sqrt{D_{\text{out}} D_{\text{in}}} \right)}{\sqrt{B}}
\end{equation}
This completes the proof. Under standard statistical learning theory, with probability at least $1 - \delta$ over the choice of the sample $S$, the expected multi-task loss is bounded by the empirical loss plus $2 \widehat{\mathcal{R}}_S(\mathcal{F}) + 3 \sqrt{\frac{\ln(2/\delta)}{2B}}$, which scales as $\mathcal{O}\left( C \sqrt{D_{\text{out}} D_{\text{in}} / B} \right)$.
This shows that keeping $C$ small and controlled (which is exactly what QCSW does to prevent quantization collapse) also stabilizes the generalization gap, whereas unconstrained model merging (like WCPR, which can have arbitrarily large coordinate-wise updates $y_{(i)}$ due to sorting outliers) exhibits unbounded Rademacher complexity and is prone to overfitting and poor generalization under noise and OOD shifts.

\subsection{Detailed Discussion on Limitations and Future Directions}
While QCSW provides a mathematically sound and empirically superior multi-dimensional alignment framework, we identify two primary limitations that present exciting avenues for future research:
\begin{enumerate}
    \item **Data-Free vs. Data-Guided Projections:** Currently, QCSW generates random orthogonal projection directions via QR decomposition of Gaussian matrices. While this ensures $100\%$ information preservation and zero-shot data-free execution, it does not account for the varying sensitivity of different parameter sub-spaces. Future work will explore data-guided or fisher-guided projections, where the projection directions $\Theta$ are aligned with the eigenvectors of the Fisher Information Matrix of each expert, further preserving critical parameter manifolds.
    \item **Scaling to Extreme Scale Parameters:** For ultra-large language models (e.g., $>70$B parameters), the input dimensionality $D_{\text{in}}$ can be extremely large (e.g., $8192$ or $16384$), making the cubic cost of QR decomposition ($\mathcal{O}(D_{\text{in}}^3)$) non-negligible. To scale to these regimes, future work will explore block-orthogonal Sliced Wasserstein distances, where weight matrices are partitioned into smaller sub-blocks before Sliced Wasserstein calibration, capping the computational complexity.
\end{enumerate}

\end{document}
